\section{Phase 5: Global Assembly Mechanisms for BLOM}
\label{sec:phase5_global_assembly}

Phase~3 and Phase~4 establish BLOM at the \emph{window level}:
for each segment index \(i\), the local window problem is well posed, and its central segment coefficient can be computed analytically.
However, this is still not a full trajectory theory.
The missing step is a \emph{global assembly mechanism}.

The core difficulty is that the \(C^{2s-2}\) continuity proven in Phase~3 is only a \emph{within-window} statement.
It does \emph{not} imply that the central segment extracted from window \(W(i,k)\) and the central segment extracted from window \(W(i+1,k)\) automatically match up to order \(2s-2\) at their shared knot.
Therefore, without an additional assembly rule, BLOM remains a collection of local variational objects rather than a mathematically complete trajectory class.

In this section, we formalize three assembly mechanisms:

\begin{itemize}
    \item \textbf{Scheme A: Shared junction states.}
    Neighboring windows share common knot-state variables and are assembled through a joint global quadratic program.
    \item \textbf{Scheme B: Overlapping consensus.}
    Each window is solved independently, and the overlapping knot-state predictions are reconciled by an explicit consensus projection.
    \item \textbf{Scheme C: Central-segment extraction.}
    Each window contributes only its central segment.
    This is the cleanest realization of the ``local coefficient map'' view, but by itself it does not automatically guarantee derivative continuity.
\end{itemize}

The three schemes have different mathematical guarantees.
The purpose of this section is not to force them into the same theorem, but to state precisely what each one does and does not guarantee.

\subsection{Common notation and the jump operator}
\label{subsec:phase5_common_notation}

We retain the canonical setting fixed in Phase~0:
\[
p:[0,\mathcal T]\to\mathbb R,\qquad
s=4,\qquad d_i=1,\qquad k=2,
\qquad
\theta=(q_1,\dots,q_{M-1},T_1,\dots,T_M)^\top,
\]
with
\[
q_0,q_1,\dots,q_M\in\mathbb R,
\qquad
T_i>0,
\qquad
\mathcal T=\sum_{i=1}^{M}T_i.
\]
For the present section, however, we keep the order \(s\ge 2\) symbolic whenever possible, and specialize to \(s=4\) when needed.

\paragraph{Knot-state vector.}
For every interior knot \(t_i\), define the lower-order knot-state vector
\begin{equation}
\label{eq:phase5_eta_i}
\eta_i
:=
\begin{bmatrix}
p^{(1)}(t_i)\\
\vdots\\
p^{(s-1)}(t_i)
\end{bmatrix}
\in \mathbb R^{s-1},
\qquad i=1,\dots,M-1.
\end{equation}
In the canonical minimum-snap case \(s=4\),
\begin{equation}
\label{eq:phase5_eta_i_s4}
\eta_i=
\chi_i=
\begin{bmatrix}
v_i\\ a_i\\ j_i
\end{bmatrix}
\in\mathbb R^3,
\end{equation}
which is exactly the local jet variable introduced in Phase~4.

\paragraph{Window-level extracted segment.}
From Phase~4, for each segment \(i\) we have a locally computed segment polynomial
\[
p_i^{\mathrm{loc}}(\tau)
=
\bigl(c_i^{\mathrm{BLOM},k}\bigr)^\top \beta_s(\tau),
\qquad \tau\in[0,T_i],
\]
where \(c_i^{\mathrm{BLOM},k}\) is the central-segment coefficient extracted from the local BLOM-Strict window.

\paragraph{Left and right knot jets induced by extracted segments.}
For every interior knot \(t_i\), define the knot jets induced by the two adjacent extracted segments:
\begin{equation}
\label{eq:phase5_eta_minus_plus}
\eta_i^{-}
:=
\begin{bmatrix}
\bigl(p_i^{\mathrm{loc}}\bigr)^{(1)}(T_i)\\
\vdots\\
\bigl(p_i^{\mathrm{loc}}\bigr)^{(s-1)}(T_i)
\end{bmatrix},
\qquad
\eta_i^{+}
:=
\begin{bmatrix}
\bigl(p_{i+1}^{\mathrm{loc}}\bigr)^{(1)}(0)\\
\vdots\\
\bigl(p_{i+1}^{\mathrm{loc}}\bigr)^{(s-1)}(0)
\end{bmatrix}.
\end{equation}
Likewise, for any order \(\ell\ge 0\), define the derivative jump at \(t_i\) by
\begin{equation}
\label{eq:phase5_jump_def}
\llbracket p^{(\ell)}\rrbracket(t_i)
:=
\bigl(p_i\bigr)^{(\ell)}(T_i)-\bigl(p_{i+1}\bigr)^{(\ell)}(0),
\qquad i=1,\dots,M-1.
\end{equation}

\paragraph{Target global continuity order.}
For BLOM as a trajectory representation, the physically essential continuity order is usually
\begin{equation}
\label{eq:phase5_target_r}
r_\star:=s-1.
\end{equation}
In the canonical minimum-snap case \(s=4\), this is \(r_\star=3\), i.e.\ continuity of position, velocity, acceleration, and jerk.
The stronger within-window continuity \(C^{2s-2}\) from Phase~3 is a local optimality consequence, but not automatically a global assembly consequence.

\subsection{A universal single-segment Hermite reconstruction map}
\label{subsec:phase5_single_segment_hermite}

All three assembly schemes will reconstruct the final global trajectory segment by segment from endpoint values and lower-order knot states.
We therefore first define the universal single-segment reconstruction operator.

Consider a single segment of duration \(T_i\) joining \(q_{i-1}\) to \(q_i\), with lower-order endpoint jets
\[
\eta_{i-1}\in\mathbb R^{s-1},
\qquad
\eta_i\in\mathbb R^{s-1}.
\]
Let
\[
\widetilde y_i
=
\begin{bmatrix}
q_{i-1}\\
D_s(T_i)\eta_{i-1}\\
q_i\\
D_s(T_i)\eta_i
\end{bmatrix}
\in\mathbb R^{2s},
\]
where
\begin{equation}
\label{eq:phase5_Ds}
D_s(T_i):=\operatorname{diag}(T_i,T_i^2,\dots,T_i^{s-1}).
\end{equation}
Let \(C_s\in\mathbb R^{2s\times 2s}\) denote the normalized degree-\((2s-1)\) Hermite interpolation matrix, and \(H_s:=C_s^{-1}\).
Let
\begin{equation}
\label{eq:phase5_lambda_2s}
\Lambda_{2s}(T_i):=\operatorname{diag}(1,T_i^{-1},\dots,T_i^{-(2s-1)}).
\end{equation}
Then the exact physical coefficient vector of the segment is
\begin{equation}
\label{eq:phase5_segment_reconstruction}
c_i
=
\mathcal H_i(q_{i-1},q_i,T_i,\eta_{i-1},\eta_i)
:=
\Lambda_{2s}(T_i)\,H_s\,\widetilde y_i.
\end{equation}
This is the same Hermite reconstruction principle already used in Phase~4, written now in general order \(s\).

\begin{proposition}[Exact segment interpolation and \(C^{s-1}\)-compatibility criterion]
\label{prop:phase5_basic_hermite}
For every segment \(i\), the coefficient vector \eqref{eq:phase5_segment_reconstruction} defines the unique polynomial of degree at most \(2s-1\) satisfying
\[
p_i(0)=q_{i-1},\qquad p_i(T_i)=q_i,
\qquad
p_i^{(\ell)}(0)=\eta_{i-1}^{(\ell)},
\qquad
p_i^{(\ell)}(T_i)=\eta_i^{(\ell)},
\quad \ell=1,\dots,s-1.
\]
Consequently, if the assembled global trajectory uses the same shared knot-state \(\eta_i\) on both sides of the interior knot \(t_i\), then
\[
\llbracket p^{(\ell)}\rrbracket(t_i)=0,
\qquad \ell=0,1,\dots,s-1.
\]
\end{proposition}

\begin{proof}
Existence and uniqueness follow from unisolvence of Hermite interpolation of degree \(2s-1\) with endpoint derivatives of orders \(0,\dots,s-1\).
The jump conclusion follows immediately, because both adjacent segments are reconstructed with the same endpoint value \(q_i\) and the same lower-order knot-state \(\eta_i\).
\end{proof}

\subsection{Scheme A: shared junction states}
\label{subsec:phase5_schemeA}

Scheme~A introduces global shared knot states as primary assembly variables.
This is the cleanest optimization-based global assembly mechanism, but it slightly weakens the purely local character of BLOM because neighboring windows are coupled through shared variables.

\subsubsection{Definition of the shared-state assembly energy}
\label{subsubsec:phase5_schemeA_definition}

For each interior knot \(i\in\{1,\dots,M-1\}\), introduce a global unknown
\[
\eta_i\in\mathbb R^{s-1}.
\]
Let
\[
\eta:=(\eta_1^\top,\dots,\eta_{M-1}^\top)^\top\in\mathbb R^{(s-1)(M-1)}.
\]
For every window \(W(r,k)\), define the local BLOM-Strict minimum energy under prescribed internal knot states:
\begin{equation}
\label{eq:phase5_local_energy_A}
\mathcal E_{r,k}^{A}(\eta;q,T)
:=
\min_{u\in\mathcal A_{r,k}^{A}(\eta;q,T)} J_{r,k}[u],
\end{equation}
where \(\mathcal A_{r,k}^{A}(\eta;q,T)\) is the Phase~3 admissible set augmented with the constraints
\begin{equation}
\label{eq:phase5_shared_state_constraints}
u^{(\ell)}(t_j)=\eta_j^{(\ell)},
\qquad
\ell=1,\dots,s-1,
\qquad
j\in \{L_r,\dots,R_r-1\}.
\end{equation}
Here \(L_r=\min W(r,k)\) and \(R_r=\max W(r,k)\), and the derivative values are interpreted at the corresponding local knots of the window.

Because the local problem is quadratic and all added constraints are affine, \(\mathcal E_{r,k}^{A}\) is a quadratic function of the subset of \(\eta\) that appears in window \(W(r,k)\):
\begin{equation}
\label{eq:phase5_local_energy_A_quad}
\mathcal E_{r,k}^{A}(\eta;q,T)
=
\frac{1}{2}\eta^\top H_{r,k}^{A}(T)\eta
-
\eta^\top g_{r,k}^{A}(q,T)
+
\gamma_{r,k}^{A}(q,T),
\end{equation}
where \(H_{r,k}^{A}(T)\) is sparse and supported only on the block indices belonging to the window.

We then define the global Scheme~A assembly problem by
\begin{equation}
\label{eq:phase5_schemeA_global_problem}
\boxed{
\min_{\eta\in\mathbb R^{(s-1)(M-1)}}
\mathcal J_A(\eta;q,T)
:=
\sum_{r=1}^{M}\mathcal E_{r,k}^{A}(\eta;q,T).
}
\end{equation}

\subsubsection{Resulting assembled trajectory}
\label{subsubsec:phase5_schemeA_reconstruction}

Given a minimizer \(\eta^\star\) of \eqref{eq:phase5_schemeA_global_problem}, define the assembled segment coefficients by the universal Hermite map
\begin{equation}
\label{eq:phase5_schemeA_ci}
c_i^{A}
=
\mathcal H_i(q_{i-1},q_i,T_i,\eta_{i-1}^\star,\eta_i^\star),
\qquad i=1,\dots,M,
\end{equation}
with the convention that the boundary jet is inherited from Phase~1 at \(i=0\) and \(i=M\).
The resulting assembled trajectory is denoted by
\[
p^{A}=(p_1^A,\dots,p_M^A).
\]

\subsubsection{Continuity theorem for Scheme A}
\label{subsubsec:phase5_schemeA_theorem}

\begin{theorem}[Scheme A yields exact global \(C^{s-1}\) continuity]
\label{thm:phase5_schemeA_continuity}
Assume that the quadratic program \eqref{eq:phase5_schemeA_global_problem} admits a unique minimizer \(\eta^\star\).
Then the assembled trajectory \(p^A\) defined by \eqref{eq:phase5_schemeA_ci} satisfies
\[
\llbracket (p^A)^{(\ell)} \rrbracket(t_i)=0,
\qquad
\ell=0,1,\dots,s-1,
\qquad
i=1,\dots,M-1.
\]
Hence \(p^A\in C^{s-1}([0,\mathcal T])\).
\end{theorem}

\begin{proof}
By construction, both segment \(i\) and segment \(i+1\) are reconstructed using the same shared knot state \(\eta_i^\star\) at the knot \(t_i\), together with the same waypoint value \(q_i\).
The claim therefore follows directly from Proposition~\ref{prop:phase5_basic_hermite}.
\end{proof}

\begin{theorem}[Scheme A existence and uniqueness under positive definiteness]
\label{thm:phase5_schemeA_existence}
Assume that the assembled Hessian
\begin{equation}
\label{eq:phase5_HA}
H_A(T):=\sum_{r=1}^{M} H_{r,k}^{A}(T)
\end{equation}
is symmetric positive definite.
Then the global Scheme~A problem \eqref{eq:phase5_schemeA_global_problem} admits a unique minimizer
\[
\eta^\star = H_A(T)^{-1}\,g_A(q,T),
\qquad
g_A(q,T):=\sum_{r=1}^{M} g_{r,k}^{A}(q,T).
\]
\end{theorem}

\begin{proof}
Under the stated assumption, \(\mathcal J_A(\eta;q,T)\) is a strictly convex quadratic function of \(\eta\).
Hence it has a unique minimizer characterized by the linear system
\[
H_A(T)\eta^\star=g_A(q,T).
\]
\end{proof}

\begin{remark}[Marked limitation for Scheme A]
\label{rem:phase5_schemeA_limitation}
A \emph{general}, assumption-free proof that \(H_A(T)\) is automatically positive definite for arbitrary \((s,k,M)\) is \textbf{not completed in this chapter}.
What is completed here is:
\begin{enumerate}
    \item the exact definition of Scheme~A,
    \item the explicit global quadratic assembly objective,
    \item the exact \(C^{s-1}\) continuity theorem,
    \item and the existence/uniqueness theorem under the explicit condition \(H_A(T)\succ 0\).
\end{enumerate}
This is the correct rigorous statement at the current stage.
\end{remark}

\subsection{Scheme B: overlapping consensus}
\label{subsec:phase5_schemeB}

Scheme~B is closest to the sliding-window philosophy:
each window is solved independently, and the overlapping knot predictions are reconciled afterwards.
Unlike Scheme~A, the independent local windows are never coupled during the local solve phase.

\subsubsection{Independent local predictions}
\label{subsubsec:phase5_schemeB_local_predictions}

For every window \(W(r,k)\), let
\[
u_{r,k}^\star
\]
be the unique BLOM-Strict local minimizer from Phase~3.
Whenever an interior knot \(t_i\) lies inside that window, define the window-specific prediction of the lower-order knot state by
\begin{equation}
\label{eq:phase5_local_eta_prediction}
\widehat\eta_i^{(r)}
:=
\begin{bmatrix}
\bigl(u_{r,k}^\star\bigr)^{(1)}(t_i)\\
\vdots\\
\bigl(u_{r,k}^\star\bigr)^{(s-1)}(t_i)
\end{bmatrix}
\in\mathbb R^{s-1}.
\end{equation}
Let
\[
\mathcal I(i):=\{\,r\in\{1,\dots,M\}: t_i \text{ is an internal knot of } W(r,k)\,\}
\]
be the set of windows that provide a prediction for knot \(i\).

\subsubsection{Consensus projection}
\label{subsubsec:phase5_schemeB_consensus}

Choose positive weights
\[
\omega_i^{(r)}>0,\qquad r\in\mathcal I(i),
\]
and define the consensus objective
\begin{equation}
\label{eq:phase5_schemeB_consensus_obj}
\mathcal J_B(\eta)
:=
\frac12\sum_{i=1}^{M-1}\sum_{r\in\mathcal I(i)}
\omega_i^{(r)}
\bigl\|\eta_i-\widehat\eta_i^{(r)}\bigr\|_2^2.
\end{equation}
Because this objective is separable in the knot index \(i\), the minimizer is given in closed form.

\begin{proposition}[Closed-form consensus state]
\label{prop:phase5_schemeB_closed_form}
The unique minimizer of \eqref{eq:phase5_schemeB_consensus_obj} is
\begin{equation}
\label{eq:phase5_schemeB_closed_form}
\eta_i^\star
=
\frac{\sum_{r\in\mathcal I(i)} \omega_i^{(r)}\,\widehat\eta_i^{(r)}}
{\sum_{r\in\mathcal I(i)} \omega_i^{(r)}},
\qquad i=1,\dots,M-1.
\end{equation}
\end{proposition}

\begin{proof}
For each fixed \(i\), the objective contribution is
\[
\mathcal J_{B,i}(\eta_i)
=
\frac12\sum_{r\in\mathcal I(i)} \omega_i^{(r)}
\|\eta_i-\widehat\eta_i^{(r)}\|_2^2.
\]
Differentiating gives
\[
\nabla_{\eta_i}\mathcal J_{B,i}(\eta_i)
=
\left(\sum_{r\in\mathcal I(i)}\omega_i^{(r)}\right)\eta_i
-
\sum_{r\in\mathcal I(i)}\omega_i^{(r)}\widehat\eta_i^{(r)}.
\]
Since all weights are positive, the coefficient matrix is a positive scalar multiple of the identity, hence invertible.
Setting the gradient to zero yields \eqref{eq:phase5_schemeB_closed_form}.
\end{proof}

\subsubsection{Reconstruction and continuity}
\label{subsubsec:phase5_schemeB_reconstruction}

Using the consensus knot states \(\eta^\star\), define the assembled segment coefficients exactly as in Scheme~A:
\begin{equation}
\label{eq:phase5_schemeB_ci}
c_i^{B}
=
\mathcal H_i(q_{i-1},q_i,T_i,\eta_{i-1}^\star,\eta_i^\star),
\qquad i=1,\dots,M.
\end{equation}
The corresponding assembled trajectory is
\[
p^B=(p_1^B,\dots,p_M^B).
\]

\begin{theorem}[Scheme B yields exact global \(C^{s-1}\) continuity]
\label{thm:phase5_schemeB_continuity}
The assembled trajectory \(p^B\) satisfies
\[
\llbracket (p^B)^{(\ell)} \rrbracket(t_i)=0,
\qquad
\ell=0,1,\dots,s-1,
\qquad
i=1,\dots,M-1.
\]
Hence \(p^B\in C^{s-1}([0,\mathcal T])\).
\end{theorem}

\begin{proof}
The proof is identical to that of Theorem~\ref{thm:phase5_schemeA_continuity}.
Both neighboring segments are reconstructed from the same consensus knot state \(\eta_i^\star\) and the same waypoint \(q_i\), so the jump vanishes up to order \(s-1\).
\end{proof}

\begin{remark}[What Scheme B guarantees and what it does not]
\label{rem:phase5_schemeB_scope}
Scheme~B guarantees exact global \(C^{s-1}\) continuity after the consensus projection.
It does \emph{not} imply that the raw independent window solutions already agree in higher derivatives.
In particular, no automatic \(C^{2s-2}\) theorem is claimed at the global level.
\end{remark}

\subsection{Scheme C: central-segment extraction}
\label{subsec:phase5_schemeC}

Scheme~C is the cleanest realization of the local-map interpretation of BLOM:
each window is solved independently, and only its central segment is kept.
No extra global variables and no consensus step are introduced.

\subsubsection{Definition}
\label{subsubsec:phase5_schemeC_definition}

For each segment \(i\), let
\[
c_i^{C}:=c_i^{\mathrm{BLOM},k}
\]
be the coefficient vector extracted directly from the local window \(W(i,k)\), as defined in Phase~4.
Define the assembled piecewise trajectory
\begin{equation}
\label{eq:phase5_schemeC_trajectory}
p^C(t):=p_i^C(t-t_{i-1}),
\qquad
t\in[t_{i-1},t_i),
\qquad
p_i^C(\tau):=(c_i^C)^\top\beta_s(\tau).
\end{equation}

\subsubsection{Exact interpolation and the guaranteed \(C^0\) property}
\label{subsubsec:phase5_schemeC_C0}

\begin{theorem}[Scheme C always preserves exact waypoint interpolation]
\label{thm:phase5_schemeC_C0}
The raw central-segment extraction trajectory \(p^C\) satisfies
\[
p^C(t_i^-)=p^C(t_i^+)=q_i,
\qquad i=1,\dots,M-1.
\]
Hence \(p^C\in C^0([0,\mathcal T])\).
\end{theorem}

\begin{proof}
By Phase~3 and Phase~4, every locally extracted central segment interpolates its two endpoint waypoints exactly.
Therefore segment \(i\) ends at \(q_i\), and segment \(i+1\) starts at \(q_i\).
Thus the value jump at every interior knot is zero.
\end{proof}

\subsubsection{Jump formulas}
\label{subsubsec:phase5_schemeC_jump_formula}

For the raw extraction trajectory, derivative continuity beyond order \(0\) is not automatic.
Define the left and right endpoint jet maps induced by the extracted segments:
\begin{equation}
\label{eq:phase5_schemeC_endpoint_maps}
\Gamma_i^{-}(q,T)
:=
\begin{bmatrix}
\bigl(p_i^C\bigr)^{(1)}(T_i)\\
\vdots\\
\bigl(p_i^C\bigr)^{(s-1)}(T_i)
\end{bmatrix},
\qquad
\Gamma_i^{+}(q,T)
:=
\begin{bmatrix}
\bigl(p_{i+1}^C\bigr)^{(1)}(0)\\
\vdots\\
\bigl(p_{i+1}^C\bigr)^{(s-1)}(0)
\end{bmatrix}.
\end{equation}
Then
\begin{equation}
\label{eq:phase5_schemeC_jump_lower}
\begin{bmatrix}
\llbracket (p^C)^{(1)}\rrbracket(t_i)\\
\vdots\\
\llbracket (p^C)^{(s-1)}\rrbracket(t_i)
\end{bmatrix}
=
\Gamma_i^{-}(q,T)-\Gamma_i^{+}(q,T).
\end{equation}
Likewise, for any \(\ell\in\{s,\dots,2s-2\}\), the higher-order jump is
\begin{equation}
\label{eq:phase5_schemeC_jump_higher}
\llbracket (p^C)^{(\ell)}\rrbracket(t_i)
=
B^{(\ell)}(T_i)c_i^C - B^{(\ell)}(0)c_{i+1}^C,
\end{equation}
where \(B^{(\ell)}(\tau)\) is the derivative row operator introduced in Phase~1.

\begin{proposition}[Sufficient condition for higher continuity in Scheme C]
\label{prop:phase5_schemeC_sufficient}
Fix an integer \(r\in\{1,\dots,s-1\}\).
If, for every interior knot \(t_i\),
\begin{equation}
\label{eq:phase5_schemeC_compatibility}
\Gamma_i^{-}(q,T)=\Gamma_i^{+}(q,T)
\quad\text{up to order }r,
\end{equation}
then the raw central extraction trajectory \(p^C\) belongs to \(C^r([0,\mathcal T])\).
\end{proposition}

\begin{proof}
Equation \eqref{eq:phase5_schemeC_jump_lower} shows that the derivative jumps of orders \(1,\dots,r\) vanish exactly under the compatibility assumption \eqref{eq:phase5_schemeC_compatibility}.
Combined with Theorem~\ref{thm:phase5_schemeC_C0}, this yields \(C^r\).
\end{proof}

\begin{remark}[Marked limitation for Scheme C]
\label{rem:phase5_schemeC_limitation}
A theorem asserting that \emph{raw} central-segment extraction automatically achieves \(C^r\) for any \(r\ge 1\) is \textbf{not completed here}, and should \textbf{not} be claimed in general.
In fact, without an additional compatibility mechanism, such a theorem is generally false:
adjacent extracted segments come from different local windows, and their endpoint jets need not coincide.
Therefore Scheme~C, in its raw form, is rigorously guaranteed only to be \(C^0\), plus whatever higher regularity can be certified by the explicit jump formulas \eqref{eq:phase5_schemeC_jump_lower}--\eqref{eq:phase5_schemeC_jump_higher}.
\end{remark}

\subsection{Comparison of the three schemes}
\label{subsec:phase5_comparison}

We now summarize the precise mathematical guarantees.

\begin{table}[t]
\centering
\caption{Mathematical guarantees of the three BLOM assembly schemes.}
\label{tab:phase5_scheme_compare}
\begin{tabular}{lccc}
\hline
Property & Scheme A & Scheme B & Scheme C \\
\hline
No extra global variables & No & Yes & Yes \\
Independent local window solves & No & Yes & Yes \\
Closed-form global assembly step & Linear system & Weighted average & None \\
Guaranteed global \(C^0\) & Yes & Yes & Yes \\
Guaranteed global \(C^{s-1}\) & Yes & Yes & No (conditional only) \\
Guaranteed global \(C^{2s-2}\) & No claim & No claim & No claim \\
Theorem completed here & Conditional & Yes & Partial \\
\hline
\end{tabular}
\end{table}

\paragraph{Interpretation.}
Scheme~A is the most theoretically coherent optimization-based assembly mechanism, but it weakens pure locality by introducing shared variables and, in full generality, still needs an external positive-definiteness condition for the assembled Hessian.
Scheme~B is the cleanest way to preserve sliding-window independence while still guaranteeing exact global \(C^{s-1}\) continuity after the projection step.
Scheme~C is the purest embodiment of the BLOM local-map philosophy, but by itself it only guarantees exact \(C^0\) continuity.

\subsection{Which scheme should be preferred?}
\label{subsec:phase5_preference}

If the main objective is a mathematically closed and practically implementable assembly mechanism with exact global \(C^{s-1}\) continuity, then Scheme~B is the safest choice:
it separates local computation from global projection and admits an explicit closed-form assembly rule.

If the main objective is to preserve the interpretation
\[
\text{``each local window defines only its central segment''},
\]
then Scheme~C is conceptually the most faithful.
However, in that case one must accept that higher-order continuity is no longer automatic and must be diagnosed explicitly through jump statistics.

Scheme~A is the most elegant in terms of variational structure, but at the present stage its complete unconditional theory is not yet as closed as that of Scheme~B.

\subsection{Canonical specialization to \texorpdfstring{$s=4,\ k=2$}{s=4,k=2}}
\label{subsec:phase5_canonical_s4k2}

In the canonical setting of this work,
\[
s=4,\qquad k=2.
\]
Hence
\[
\eta_i=\chi_i=
\begin{bmatrix}
v_i\\ a_i\\ j_i
\end{bmatrix}\in\mathbb R^3.
\]
The target global continuity order is
\[
r_\star=s-1=3.
\]
Accordingly:
\begin{itemize}
    \item Scheme~A and Scheme~B both guarantee
    \[
    \llbracket p\rrbracket(t_i)=
    \llbracket p'\rrbracket(t_i)=
    \llbracket p''\rrbracket(t_i)=
    \llbracket p^{(3)}\rrbracket(t_i)=0,
    \]
    provided their assembly problems are solved as defined above;
    \item Scheme~C guarantees only
    \[
    \llbracket p\rrbracket(t_i)=0
    \]
    unconditionally, while
    \[
    \llbracket p'\rrbracket(t_i),\ 
    \llbracket p''\rrbracket(t_i),\ 
    \llbracket p^{(3)}\rrbracket(t_i)
    \]
    must be checked explicitly through \eqref{eq:phase5_schemeC_jump_lower};
    \item none of the three schemes, in the present chapter, is claimed to guarantee automatic global \(C^6\) continuity.
\end{itemize}

\subsection{Numerical consequence and Phase~5 validation target}
\label{subsec:phase5_validation_target}

The central numerical diagnostic of Phase~5 is the jump statistic
\[
\llbracket p^{(\ell)}\rrbracket(t_i),
\qquad
\ell=0,\dots,2s-2,
\qquad
i=1,\dots,M-1.
\]
For Scheme~A and Scheme~B, the theory predicts exact cancellation for \(\ell=0,\dots,s-1\).
For Scheme~C, only \(\ell=0\) is guaranteed exactly, while \(\ell\ge 1\) should generally be treated as a diagnostic quantity.
This makes the boundary-jump checker the primary numerical witness for deciding which global assembly mechanism is appropriate for the intended BLOM application.